\documentclass[12pt]{article}
\usepackage[utf8]{inputenc}
\usepackage[T1]{fontenc}
\usepackage[margin=1in]{geometry}
\usepackage{setspace}
\usepackage{xcolor}
\usepackage[colorlinks=true, urlcolor=blue!60!black, linkcolor=blue!60!black]{hyperref}
\setlength{\parindent}{2em}
\setlength{\parskip}{1.25em}
\renewcommand{\baselinestretch}{1.0}
\begin{document}
\begin{center}
{\LARGE \textit{Garcia v. San Antonio Metropolitan Transit Authority}, 469 U.S. 528 (1985)\par}
\vspace{0.8em}
{\large Case Brief --- Four-Page Edition\par}
\end{center}
\vspace{1em}

\section{Memory Jogger}

Overruling \textit{National League of Cities v. Usery}, the Court held 5--4 that a city transit authority must obey the Fair Labor Standards Act's wage and overtime rules: ``traditional governmental functions'' immunity is unworkable, and the States are protected mainly by the political process (\href{https://supreme.justia.com/cases/federal/us/469/528/}{Justia}).

\section{Detailed Case Facts}

San Antonio's transit system became public in 1959, and in 1978 its facilities passed to the San Antonio Metropolitan Transit Authority (SAMTA). It received over \$51 million in federal Urban Mass Transportation Act grants between 1970 and 1980. Congress had extended the FLSA to public mass transit and, in 1974, to almost all state and local employees. After \textit{National League of Cities} (1976) barred enforcing the FLSA against States ``in areas of traditional governmental functions,'' the transit system told employees it was no longer bound by the overtime rules.

\section{Procedural History}

In 1979 the Labor Department found SAMTA not immune. SAMTA sued the Secretary of Labor for a declaratory judgment; the Secretary counterclaimed, and employee Joe Garcia intervened. The district court held transit a traditional function. After \textit{United Transportation Union v. Long Island R. Co.} (1982) the Court vacated and remanded; the district court adhered, comparing the inquiry to identifying pornography. The Supreme Court heard argument in 1984 and ordered reargument on whether to reconsider \textit{National League of Cities}.

\section{Judicial Votes}

Reversed and remanded, \textbf{5--4} (decided February 19, 1985).
\begin{itemize}
\item \textbf{Majority:} Justice Blackmun, joined by Justices Brennan, White, Marshall, and Stevens.
\item \textbf{Dissents:} Justice Powell, joined by Chief Justice Burger and Justices Rehnquist and O'Connor; Justice Rehnquist; Justice O'Connor, joined by Justices Powell and Rehnquist.
\end{itemize}

\section{Holding}

\textit{National League of Cities} is overruled. Defining state immunity by identifying traditional governmental functions is unworkable in practice and unsound in principle. Because the Constitution protects state sovereignty chiefly through procedural safeguards built into the national political process, applying the FLSA's wage and overtime rules to SAMTA violated no affirmative limit on the commerce power. Blackmun had supplied the fifth vote in \textit{National League of Cities}.

\section{Analysis of Opinions}

\textbf{Majority (Blackmun).} Eight years of lower-court cases produced irreconcilable results, granting immunity for ambulances, licensing, airports, and solid waste but denying it for traffic regulation and mental-health facilities. The Court's earlier attempt to separate ``governmental'' from ``proprietary'' functions in tax immunity had likewise failed. Historical tests freeze state governance; nonhistorical tests invite judges to pick favored services; and both penalize States acting as laboratories of experiment. Relying on \textit{The Federalist} and Herbert Wechsler's \textit{Political Safeguards of Federalism}, Blackmun found the States' protection in the structure of the national government---state control of House qualifications, the Electoral College, and equal Senate representation. The process worked: federal grants grew from \$7 billion to \$96 billion, and States won exemptions from the NLRA, OSHA, and ERISA. SAMTA faced the same rules as private employers and received tens of millions in federal aid.

\textbf{Powell (dissenting).} The Court abandoned stare decisis after recent reaffirmances; \textit{National League of Cities} set a workable balancing test; making Congress the judge of its own power defies \textit{Marbury}; and the Seventeenth Amendment, national media, and lobbying had eroded Wechsler's premises.

\textbf{Rehnquist (dissenting).} SAMTA should win under any dissent's formulation, and the principle will in time again command a majority.

\textbf{O'Connor (dissenting).} Under \textit{McCulloch}, means must conform to the letter and spirit of the Constitution, including the Tenth Amendment's ``double security.''

\textbf{Later.} \textit{New York}, \textit{Printz}, and \textit{Murphy} bar commandeering; \textit{Reno v. Condon} reaffirmed generally applicable rules; \textit{Alden v. Maine} bars private FLSA damages suits against States in their own courts.

\section{Examples: Future Applications}

\begin{enumerate}
\item \textbf{Same-Side (Valid): Bus Electrification.} A rule requiring all public and private bus operators to electrify half their fleets regulates market participants generally.
\item \textbf{Same-Side (Valid): Shift Caps.} A generally applicable cap on emergency workers' shifts binds municipal police and firefighters, who must seek relief in Congress.
\item \textbf{Opposite-Side (Invalid): DMV Enforcers.} Ordering state DMVs to audit private operators commandeers state officials.
\item \textbf{Opposite-Side (Barred): Private Suits in State Court.} The FLSA binds the State, but \textit{Alden} bars private overtime suits against it in its own courts.
\item \textbf{Fence-Sitter: State Judicial Staff.} Pay and hours rules only for state-court clerks and legislative staff are not generally applicable and reach core state structure.
\end{enumerate}

\section{Critique}

\textbf{Eroded safeguards and an undefined standard.} Critics echo Powell that the Seventeenth Amendment and nationalized politics align legislators with ideological coalitions, not state autonomy. Blackmun allowed intervention for possible failings in the political process without defining one.

\textbf{Formalism versus burden.} Modern doctrine asks \emph{how} Congress regulates, not \emph{what} it regulates, though general labor mandates can burden state budgets far more than minor commandeered duties.

\section{Key Quotations}

\begin{quote}
``We therefore now reject, as unsound in principle and unworkable in practice, a rule of state immunity from federal regulation that turns on a judicial appraisal of whether a particular governmental function is `integral' or `traditional.' Any such rule leads to inconsistent results at the same time that it disserves principles of democratic self-governance, and it breeds inconsistency precisely because it is divorced from those principles.''
\end{quote}
\noindent --- Justice Blackmun, for the Court, 469 U.S. at 546--47.

\begin{quote}
``Nonetheless, against this background, we are convinced that the fundamental limitation that the constitutional scheme imposes on the Commerce Clause to protect the `States as States' is one of process, rather than one of result. Any substantive restraint on the exercise of Commerce Clause powers must find its justification in the procedural nature of this basic limitation, and it must be tailored to compensate for possible failings in the national political process, rather than to dictate a `sacred province of state autonomy.'\,''
\end{quote}
\noindent --- Justice Blackmun, for the Court, 469 U.S. at 554.

\begin{quote}
``State sovereign interests, then, are more properly protected by procedural safeguards inherent in the structure of the federal system than by judicially created limitations on federal power.''
\end{quote}
\noindent --- Justice Blackmun, for the Court, 469 U.S. at 552.

\end{document}
